\section{System Architecture and Marker Selection Rationale}
\label{rationale}
Although it shares the same underlying concept as the previous study \cite{Owais2026}, this method is further refined to achieve a significantly higher sub-pixel precision level by replacing the ArUco fiducial marker with a chessboard marker. Based on its geometric characteristics, the chessboard marker offers vastly superior performance for sub-pixel corner detection requirements compared to ArUco \cite{Triyono2024, Dan2022}. Through sub-pixel corner detection using algorithms such as EDLines (Edge Drawing Lines), the chessboard marker can yield a highly precise average reprojection error ranging from 0.044 to 0.05 pixels \cite{Dan2022}. On the other hand, while ArUco's primary advantage lies in its functionality as a unique identifier that is highly robust against environmental noise, its spatial corner detection sharpness does not match the clear line intersections of a chessboard marker. 

As a hybrid solution, the ChArUco algorithm is frequently deployed because it combines the geometric sharpness of a chessboard with the robust identification of ArUco to handle visual occlusion and extreme environmental conditions, achieving a pose detection success rate of 97.4\cite{Hu2019}. However, a ChArUco implementation is considered overkill within the scope of this study, as its unique code identification feature becomes redundant in an ideal, controlled 3D printer workspace that is free from occlusion issues. 

The accurate corner detection results obtained from this chessboard are subsequently utilized in the pose estimation process to secure precise positional readouts. In this study, the algorithm selected for pose estimation is the Iterative SolvePnP (Perspective-n-Point) \cite{opencv_solvepnp}. The selection of the Iterative SolvePnP algorithm—the default type in OpenCV—is based on its performance, which is deemed sufficient for pose estimation when coupled with a 3D model featuring highly accurate structural features like a chessboard \cite{Zhou2025}. In summary, a precise corner detection approach yields a sharper pose estimation accuracy compared to using standard ArUco markers. Through this approach, the measurement error in kinematic tracking is expected to be suppressed to a level significantly lower than that of the prior study, establishing a calibration solution that not only rivals the accuracy of industrial optical instruments but also incurs an exceptionally low implementation cost.